
\section{The Yeh--Petroff dam break problem}
Yeh and Petroff of the University of Washington conducted experimental dam break with an obstruction column, as shown in Figure~\ref{fig:YP}.
The Yeh--Petroff dam break problem was further simulated by Gomez-Gesteira and Dalrymple~\cite{GD2004} and Silvester and Cleary~\cite{SC2006}. A similar problem was studied by Arnason et al.~\cite{APY2009}.
We shall compare our \anuga{} solution to the experimental data of Yeh and Petroff.

\begin{figure}
\begin{center}
\includegraphics[width=0.9\textwidth]{Yeh_Petroff.png}
\end{center}
\caption{The initial condition and geometry of the Yeh--Petroff dam break
problem, drawn from the values the case uses: a 1.60 by 0.61~m flume, a gate
at $x = 0.4$~m holding 0.30~m of water against a 0.01~m wet bed, and a
0.12~m square column at $x = 0.9$~m on the centre of the channel.}
\label{fig:YP}
\end{figure}




\subsection{Results}

We should see broad agreement between the experimental data and the numerical solution. Some disagreement is expected because the flows in this experimental situation show fairly rapid spatial and temporal variations, which are to some extent outside the domain of the shallow water approximation.

\begin{figure}
\begin{center}
\includegraphics[width=0.9\textwidth]{force.png}
\end{center}
\caption{Force on the column: \anuga{} against the measurements of Yeh and Petroff.}
\end{figure}
%
%
%\begin{figure}
%\begin{center}
%\includegraphics[width=0.9\textwidth]{xmom_plot.png}
%\end{center}
%\caption{Xmomentum results}
%\end{figure}
%
%
%\begin{figure}
%\begin{center}
%\includegraphics[width=0.9\textwidth]{xvel_plot.png}
%\end{center}
%\caption{Xvelocity results}
%\end{figure}


\endinput
